% ============================================================
% 1. FIRST QUESTION
% ============================================================
\section{Visualisation, Chi-Square and t-Test}\label{sec:q1}

% ------------------------------------------------------------
% a)
% ------------------------------------------------------------
\subsection{Data Import and Variable Types}\label{sec:q1a}

\task{Import the data set to R assigning the type of each variable correctly.}

We imported the data set \texttt{bodyPerformance.csv} and renamed its columns to short R names without spaces or symbols (e.g.\ \texttt{body fat\_\%} becomes \texttt{body\_fat\_pct}).

Then we gave each variable the type that matches its measurement scale:

\begin{itemize}
  \item \textbf{Nominal categorical:} \texttt{gender} is a factor, and we labeled its codes \texttt{F} and \texttt{M} as \texttt{Female} and \texttt{Male} respectively.
  \item \textbf{Ordinal categorical:} \texttt{class} is an ordered factor with the levels $D < C < B < A$, from poor to excellent performance.
  \item \textbf{Numerical:} the other ten variables (\texttt{age}, \texttt{height\_cm}, \texttt{weight\_kg}, \texttt{body\_fat\_pct}, \texttt{diastolic}, \texttt{systolic}, \texttt{grip\_force}, \texttt{sit\_bend\_cm}, \texttt{sit\_ups} and \texttt{broad\_jump\_cm}) are numeric.
    \texttt{age} could be stored as an integer, but we kept it numeric for simplicity.
\end{itemize}

The result has 13,393 observations of 12 variables with no missing values (R output, Appendix~\ref{app:code}).
The name, type and levels of each variable are listed in \autoref{tab:q1a-variables}.

% ------------------------------------------------------------
% b)
% ------------------------------------------------------------
\subsection{Descriptive Statistics}\label{sec:q1b}

\task{Summarize the numerical variables and analyze their distributions.}

The descriptive statistics of the cleaned numerical variables are presented in \autoref{tab:q1b-summary-clean}.
\paragraph{Outliers.} 
\autoref{tab:q1b-summary} summarizes the ten numerical variables of the original data set, and their boxplots (\autoref{fig:q1b-boxplots}) show values beyond $1.5 \times \mathrm{IQR}$ from the quartiles in every variable except \texttt{age} and \texttt{sit\_ups} (\autoref{tab:q1b-outliers}).
We first examined the descriptive statistics (\autoref{tab:q1b-summary}) and boxplots (\autoref{fig:q1b-boxplots}) to identify potential outliers (\autoref{tab:q1b-outliers}).

We then further evaluated the recorded values using established health reference ranges (\autoref{tab:q1b-reference}) and the characteristics of each measurement.
Based on this analysis, we excluded observations with impossible or implausible values, including extreme blood pressure readings, body fat percentages, BMI values and invalid performance test measurements.
The exclusion criteria and the number of affected observations are presented in \autoref{tab:q1b-removed}.
In total, 651 observations (4.9\%) were removed, leaving 12,742 in \texttt{bodyPerformance\_clean.csv} for subsequent analyses.

\paragraph{Distribution and normality.}
The distributions of the cleaned numerical variables were examined using histograms (\autoref{fig:q1b-histograms}) and Q-Q plots (\autoref{fig:q1b-qqplots}).
The findings are summarized in \autoref{tab:q1b-distributions}.
The Q-Q plots (\autoref{fig:q1b-qqplots}) show deviations from normality, particularly in the tails, and the Shapiro-Wilk tests rejected normality for all variables ($p < 0.05$, \autoref{tab:q1b-shapiro}).
However, based on the histograms (\autoref{fig:q1b-histograms}), several variables appear approximately bell-shaped, suggesting that a normal approximation may be reasonable for further statistical analysis.

\begin{table}[htbp]
\centering
\caption{Distribution shapes of the cleaned numerical variables.}
\label{tab:q1b-distributions}
\small
\begin{tabular}{ll}
\toprule
Variable & Distribution shape \\
\midrule
Age          & Right-skewed \\
Height       & Approximately normal \\
Weight       & Slightly right-skewed \\
Body fat     & Approximately normal \\
Diastolic    & Approximately normal \\
Systolic     & Approximately normal \\
Grip force   & Bimodal \\
Sit and bend & Left-skewed \\
Sit-ups      & Left-skewed \\
Broad jump   & Slightly left-skewed \\
\bottomrule
\end{tabular}
\end{table}
\xfcomment{TODO: describe the shape of each variable with the histograms of the cleaned data (\autoref{fig:q1b-histograms}) and compare it with the mean and median of \autoref{tab:q1b-summary-clean}.}
\xfcomment{TODO: check whether the two peaks of \texttt{grip\_force} (around 25 and 43) correspond to the two values of \texttt{gender} (Section~\ref{sec:q1g}).}
\xfcomment{TODO: comment on the Q-Q plots of the cleaned data (\autoref{fig:q1b-qqplots}): which variables follow the line and which deviate in the tails.}
\xfcomment{TODO: report the Shapiro-Wilk test on a random sample of 5,000 cleaned observations (\autoref{tab:q1b-shapiro}) and explain why, with this sample size, it rejects normality even for the variables that look normal in \autoref{fig:q1b-histograms}.}

% ------------------------------------------------------------
% c)
% ------------------------------------------------------------
\subsection{Blood Pressure by Gender}\label{sec:q1c}

\task{Consider the variables of diastolic and systolic blood pressure.
Eliminate the zero values and check their distributions for the gender groups using proper graph(s).
Interpret your findings.}

The zero blood pressure values were removed during data cleaning (Section~\ref{sec:q1b}).
\autoref{tab:q1c-summary} summarizes the systolic and diastolic blood pressure measurements by gender.

\paragraph{Systolic blood pressure.}
Males have higher mean and median systolic blood pressure than females (\autoref{tab:q1c-summary}).
The histograms (\autoref{fig:q1c-histograms}) show an approximately normal distribution for males and a right-skewed distribution for females. 
The right-skewed female systolic distribution indicates that most female observations are concentrated at relatively lower blood pressure values, with fewer observations extending toward higher values. 
The standard deviations are similar for both genders, indicating comparable variability.

\paragraph{Diastolic blood pressure.}
Males also have higher mean and median diastolic blood pressure than females.
Both distributions are unimodal (males approximately normal, females right-skewed).
The standard deviations are nearly identical, indicating similar variability between genders.

\paragraph{Boxplots.}
The boxplots (\autoref{fig:q1c-boxplots}) confirm higher median blood pressure values among males for both measurements.
Several outliers are present, mostly at higher blood pressure values.

\xfcomment{Discuss the differences in the distributions of \texttt{diastolic} and \texttt{systolic} between the two levels of \texttt{gender}, based on your graphs.}

% ------------------------------------------------------------
% d)
% ------------------------------------------------------------
\subsection{Confidence Intervals for Systolic Blood Pressure}\label{sec:q1d}

\task{Construct 95\% confidence intervals for the mean systolic blood pressure of two gender groups and interpret them.}
\paragraph{Confidence intervals.}
The estimated means and 95\% confidence intervals for both genders are presented in \autoref{tab:q1d-t-test}.

\paragraph{Hypothesis testing.}
The true population mean systolic blood pressure is estimated to lie between the lower and upper confidence interval bounds for each gender.
The p-values are below 0.05 for both one-sample t-tests.
We therefore reject the null hypothesis that the population mean systolic blood pressure equals 120 mmHg for both genders.
The confidence intervals do not overlap, suggesting that males have a higher population mean systolic blood pressure than females, consistent with Section~\ref{sec:q1c}.


% ------------------------------------------------------------
% e)
% ------------------------------------------------------------
\subsection{Age Groups}\label{sec:q1e}

\task{Categorize the variable age into four categories cutting from the values 30, 40 and 50 and save it as a new variable into your data frame.}
We categorized age into four groups using the cut-off values 30, 40 and 50, and saved the result as the new variable \texttt{age\_cat}.
The frequencies and proportions are shown in \autoref{tab:q1e-age-groups} and \autoref{fig:q1e-age-groups}.

The largest group consists of individuals younger than 30 (43.0\%), while the smallest group is aged 40--49 (14.0\%).
The remaining individuals are distributed between ages 30--39 (20.0\%) and 50 or older (22.9\%).
\xfcomment{TODO: comment on the size of each level of \texttt{age\_cat} (e.g.\ the largest and smallest groups).}

% ------------------------------------------------------------
% f)
% ------------------------------------------------------------
\subsection{Association of Age Groups with Class and Gender}\label{sec:q1f}

\task{Analyze the association first between age groups and class and then between age groups and gender by using proper statistical method.}

We used chi-square tests of independence to examine the associations between \texttt{age\_cat} and \texttt{class}, and between \texttt{age\_cat} and \texttt{gender}.
The contingency tables are presented in \autoref{tab:q1f-age-class} and \autoref{tab:q1f-age-gender}, and the test results in \autoref{tab:q1f-chisq}.

\paragraph{Age group and performance class.}
The chi-square test indicates a statistically significant association between age group and performance class ($\chi^2(9) = 110.77$, $p < 10^{-16}$). We reject the null hypothesis of independence.
The proportions (\autoref{fig:q1f-proportions}) show that older age groups have higher proportions of individuals in lower performance classes (C and D), while younger age groups have higher proportions in classes A and B.

\paragraph{Age group and gender.}
The chi-square test also indicates a statistically significant association between age group and gender ($\chi^2(3) = 188.42$, $p < 10^{-16}$). We reject the null hypothesis of independence.
Males represent the majority in every age group, but their proportion is highest among individuals aged 30--39.
The proportion of females is highest in the 50+ age group (\autoref{fig:q1f-proportions}).
\xfcomment{TODO: state the hypotheses, check the assumption (expected counts $\geq 5$), report $\chi^2$, df and p-value, and conclude on the independence of \texttt{age\_cat} with \texttt{class} and with \texttt{gender}.}

% ------------------------------------------------------------
% g)
% ------------------------------------------------------------
\subsection{Physical Performance by Gender}\label{sec:q1g}

\task{Test whether physical performance test results of the individuals are affected from the gender.
Check the assumptions and visualize the relationship between these two characteristics.}

We used two-sample t-tests to compare physical performance results between genders.
The Shapiro-Wilk tests rejected normality for all variables (\autoref{tab:q1g-shapiro}). However, the large sample sizes make the t-tests reasonably robust to deviations from normality.
Although the Shapiro-Wilk tests rejected normality for all variables (\autoref{tab:q1g-shapiro}), the distributions were identified as approximately normal based on the analysis in Section~\ref{sec:q1b}.
The F-tests indicated unequal variances for grip strength, sit-and-bend flexibility and sit-ups, so Welch's t-test was used.
For broad jump, Student's t-test was applied (\autoref{tab:q1g-f-test}).

\paragraph{Statistical results.}
All four t-tests showed statistically significant differences between genders ($p < 0.001$), so the null hypothesis of equal means was rejected for each variable (\autoref{tab:q1g-t-test}).
Males achieved higher mean results in grip strength, sit-ups and broad jump, while females performed better in the sit-and-bend flexibility test. The mean differences were approximately 17.61 kg for grip strength, 13.85 repetitions for sit-ups, 57.91 cm for broad jump and 5.90 cm for flexibility.
The density plots (\autoref{fig:q1g-distributions}) show clear differences between genders, particularly in grip strength and broad jump, where the distributions overlap relatively little.
The boxplots (\autoref{fig:q1g-boxplots}) confirm higher median values for males in grip strength, sit-ups and broad jump, and for females in flexibility.
\xfcomment{Summarize the assumption checks and discuss whether \texttt{grip\_force}, \texttt{sit\_bend\_cm}, \texttt{sit\_ups} and \texttt{broad\_jump\_cm} differ significantly by \texttt{gender}.}
